\section{Related Work}

\paragraph{Parameter-efficient and prompt-based adaptation.}
Prefix tuning prepends learned continuous states throughout the Transformer~\citep{li2021prefix}, while prompt tuning optimizes input embeddings and becomes competitive as model scale grows~\citep{lester2021prompt}.  P-Tuning v2 extends deep prompts across scales and tasks~\citep{liu2022ptuningv2}; SPoT transfers prompts across tasks~\citep{vu2022spot}; residual prompt tuning improves optimization through a reparameterized residual path~\citep{razdaibiedina2023residual}; and InfoPrompt regularizes soft prompts using mutual-information objectives~\citep{wu2023infoprompt}.  LoRA instead learns low-rank weight updates~\citep{hu2022lora}.  Our setting is stricter: the model weights are frozen, only eight input embeddings are trainable, and the target is behavior induced by a long textual skill.

\paragraph{Context and prompt compression.}
Gist tokens learn to compress prompts into a small set of activations~\citep{mu2023gist}; AutoCompressors recursively summarize long contexts into soft memory~\citep{chevalier2023autocompressors}.  LLMLingua and its successors compress prompts by removing or reordering less informative discrete tokens~\citep{jiang2023llmlingua,jiang2024longllmlingua,pan2024llmlingua2}, while MEND learns a compact encoding for demonstrations~\citep{li2024mend}.  These methods primarily compress the input or demonstrations themselves.  \softskill instead learns from successful agent trajectories~\citep{tao2026softskill}.  On-policy context distillation optimizes a weight-updated student on its own trajectories against a privileged-context teacher~\citep{ye2026opcd}; this directly addresses the state mismatch that our frozen-prefix experiments expose, but changes many more parameters.  \method is an off-policy, parameter-efficient counterpart: it first identifies the hard skill's behavioral delta on verified trajectories.

\paragraph{Distillation and token selection.}
Knowledge distillation transfers teacher distributions to a student~\citep{hinton2015distilling}, with sequence-level KD using generated targets~\citep{kim2016sequence}.  For generative models, reverse-KL and on-policy formulations can better match the student's sampling distribution~\citep{gu2024minillm,agarwal2024gkd}.  Selective KD has been studied at the example and token levels~\citep{wang2021selective,tavor2026rethinking}; very recent work jointly adapts token selection and privileged teacher context to student capacity~\citep{yang2026usd}.  KFD is especially close in scoring tokens by divergence between a teacher and a related base model~\citep{yuce2026kfd}; Delta-KD similarly transfers the distributional shift from a raw to a fine-tuned teacher~\citep{cao2025delta}.  The distinction is the object being compressed: our teacher and reference are the \emph{same frozen agent} under hard-skill and no-skill contexts, and the only student parameters are an input prefix.  We further pair selective learning with an explicit clean-behavior preservation set and evaluate long-horizon executable generation rather than next-token accuracy alone.

\paragraph{Agent skills and prompt optimization.}
ReAct interleaves reasoning and action for language agents~\citep{yao2023react}, and AgentTuning studies instruction tuning over agent trajectories~\citep{zeng2024agenttuning}.  Natural-language prompts can be optimized using textual gradients or model-generated search directions~\citep{pryzant2023protegi,yang2024opro,yuksekgonul2024textgrad}; SkillOpt discovers reusable textual skills from successful rollouts~\citep{guo2026skillopt}.  Our work begins after a textual skill and successful trajectories exist and asks what portion of their induced behavior an extremely small continuous prefix can retain.
